% !TeX root = ../main.tex

\chapter{緒論}
\label{chap:introduction}

\section{研究背景}
\label{sec:background}

表皮內神經纖維 (Intraepidermal Nerve Fiber, IENF) 為末梢感覺神經由真皮穿越表皮／真皮交界並伸入表皮層之無髓鞘小纖維，其密度 (Intraepidermal Nerve Fiber Density, IENFD) 在近二十年間已被多項國際臨床指引認可為診斷\textbf{小纖維神經病變} (Small Fiber Neuropathy, SFN) 之黃金標準生物指標。SFN 常見於糖尿病周邊神經病變、化療誘發性神經病變、遺傳性感覺神經病變等臨床族群；由於傳統之神經傳導檢查 (Nerve Conduction Study, NCS) 僅能反映有髓鞘之粗纖維功能，無法偵測小纖維的受損，因此以皮膚切片量測 IENFD 在臨床上具有無可取代之角色。

取得 IENFD 的標準流程係於病人小腿遠端以 3 mm 環鑽取下皮膚樣本，經 PGP 9.5 免疫組織化學染色後，製成 50 $\mu m$ 之垂直切片。病理醫師或技術員於顯微鏡或影像上以人工方式計算\textbf{「跨越表皮／真皮交界 (dermal–epidermal junction, DEJ) 之神經纖維數量」}，再除以該切片之表皮長度，以每毫米之纖維數 (fibers/mm) 呈現 IENFD。該計數遵循 European Federation of Neurological Societies (EFNS) 所提之規則：一條纖維無論於表皮內如何分支，皆僅計為一次；僅計算由真皮穿越 DEJ 進入表皮之纖維，而不計算完全位於表皮內或真皮內之片段。

\section{研究動機與問題描述}
\label{sec:motivation}

雖然 IENFD 已在臨床廣泛使用，但人工計數仍面臨三項主要問題：

\begin{enumerate}
  \item \textbf{高度人力成本}：單一病人切片通常包含數十至上百條潛在纖維，計數耗時且易因疲勞產生遺漏或重複。
  \item \textbf{判讀者間變異}：不同技術員對於「一條纖維」的認定（尤其在纖維細小、折疊、部分被染色遮蔽時）缺乏一致準則，造成不同實驗室、不同時間之結果不可直接比較。
  \item \textbf{缺乏客觀中介產物}：人工計數僅產生「數字」，不產生可重複稽核的纖維拓樸結構，使得後續演算法比較與結果重現困難。
\end{enumerate}

近年已有研究嘗試以深度學習 (以 U-Net 為主) 對 PGP 9.5 染色影像進行語意分割 (semantic segmentation)，期望以\textbf{分割後再骨架化}的方式自動化 IENFD 量測。然而實務上存在一個結構性困難：\textbf{染色強度不均與切片厚度所造成的訊雜比波動，常使同一條生物學上連續之纖維於分割結果中被拆解成多個彼此不連通之「碎片」}。典型的失效情境包括：纖維穿越 DEJ 時由於表皮內染色較弱而中斷、纖維在表皮內變細至接近像素尺度、以及鄰近纖維交錯使局部訊號飽和。若直接將此破碎分割圖骨架化並計數跨越事件，將因纖維中斷而漏計大量合理之跨越，亦會因相鄰碎片的短小分支產生誤計。

欲自破碎之分割結果重建出可靠之纖維拓樸結構，須解決一個本質上是圖上之「連結 (linking)」問題：\textbf{給定一組像素空間中的連通元件，如何在不引入生物不合理連線之前提下，將原屬同一纖維之元件重新接續？} 單純以元件中心點之歐式距離做 $k$-近鄰連結（如 DBSCAN 類方法），會忽略影像中實際存在的連續訊號，在纖維密集區域產生橫向錯接；單純以最小生成樹 (MST) 串接所有元件則會產生過度連接之樹狀結構。

\section{研究目標與貢獻}
\label{sec:contributions}

本研究之核心目標為：\textbf{在不依賴額外人工干預之前提下，將由語意分割所得之破碎標註重建為與影像成本一致、可直接用於有效跨越計數之纖維拓樸結構。} 為達此目標，本文提出 \emph{標註擴張連結演算法} (Annotation-Grow Linker, AGL)。本論文之主要貢獻如下：

\begin{itemize}
  \item \textbf{基於影像之連結準則。} 不同於純粹幾何之連結方法，本文將連結問題形式化為一個建立於\textbf{像素成本圖}上的\emph{多源最短路徑}問題。成本圖由綠色通道經 CLAHE 與 Sato 血管性濾波而得，使「沿纖維行進」之成本遠低於「橫跨背景」之成本，從而使連結路徑自然遵循影像之纖維線索。

  \item \textbf{元件層圖與裁剪—MST 重建。} 以各連通元件為 Dijkstra 之獨立來源同步擴張，並於擴張場中搜尋任意兩元件之最小成本相遇點，將其萃取為元件層圖之邊；再以成本閾值裁剪異常邊後求最小生成森林，確保每個真實纖維僅保留必要之連結、避免過度連線。

  \item \textbf{像素級拓樸與跨越判定。} 將 MST 所回溯之跨元件路徑與原標註合併並骨架化，得到可供後續分段 (segment) 與跨越判定使用之像素級圖結構；在此結構上結合表皮遮罩執行分段偵測、區域標記與有效跨越計數，形成一個端到端的完整流程。

  \item \textbf{模組化之開源實作。} 演算法以 Python 實作並拆分為 \texttt{cost\_map}、\texttt{dijkstra}、\texttt{graph\_builder}、\texttt{skeleton} 等獨立模組，各模組皆可被其他連結策略重用，便於後續研究與臨床驗證。
\end{itemize}

\section{論文組織}
\label{sec:organization}

本論文後續章節組織如下。第~\ref{chap:related}~章回顧與本研究相關之醫學影像分析、血管性濾波、語意分割、以及圖論式纖維追蹤文獻，並討論既有連結策略於 IENF 重建任務之限制。第~\ref{chap:method}~章詳細介紹本文提出之 AGL 演算法，包括成本圖建構、多源 Dijkstra 擴張、元件層圖之建立與裁剪、最小生成森林、以及像素級拓樸與跨越分析之完整流程。第~\ref{chap:experiment}~章說明實驗資料集、評估指標、參數設定，並將 AGL 與既有基準方法進行定性與定量之比較，亦透過消融實驗檢視各模組之貢獻。第~\ref{chap:conclusion}~章總結本研究之發現，討論方法之限制，並指出可延伸之未來研究方向。
